% concatenated from 2025-09-15-BNW_matrices.tex
\documentclass{article}
\usepackage{graphicx} % Required for inserting images
\usepackage{a4wide} % Required for inserting images
\usepackage{amsmath, amssymb} % Required for inserting images
\usepackage{amsthm,xcolor,hyperref}
\usepackage{paralist} % very flexible & customisable lists (eg. enumerate/itemize, etc.)
\usepackage[capitalize]{cleveref}
\usepackage{amsthm, comment}

\usepackage[linesnumbered,ruled,vlined]{algorithm2e}
\SetAlgoCaptionSeparator{}
\renewcommand{\algorithmcfname}{}
\renewcommand{\thealgocf}{}

\def\alph{\tilde{\Sigma}}
\def\Id{\textsf{Id}}
\def\A{\mathcal{A}}
\let\calA = \A
\def\O{\mathcal{O}}
\def\calR{\mathcal{R}}
\def\calS{\mathcal{S}}
\def\T{\mathcal{T}}
\def\U{\mathfrak{U}}
\def\tU{\tilde{\mathfrak{U}}}
\def\tH{\tilde{H}}
\def\R{\mathfrak{R}}
\def\tM{\mathfrak{M}}
\def\tR{\tilde{\mathfrak{R}}}
\def\tS{\tilde{\mathfrak{S}}}
\def\tO{\tilde{\mathcal{O}}}
\def\tP{\tilde{P}}

\def\AA{\mathbb{A}}
\def\C{\mathbb{C}}
\def\N{\mathbb{N}}
\def\Q{\mathbb{Q}}
\def\Z{\mathbb{Z}}
\let\ZZ = \Z
\def\GL{\textsf{GL}}
\def\M{\textsf{M}}
\def\DC{\textsf{DC}}
\def\EC{\textsf{EC}}
\def\WP{\textsf{WP}}
\def\S{\textsf{QuickWP}}
\def\False{\textsf{False}}
\def\True{\textsf{True}}
\def\Ord{\textsf{Ord}}
\let\epsilon = \varepsilon
\let\phi = \varphi
\def\lab{\textsf{lab}}
\def\mat{\textsf{M}}
\def\inv{^{-1}}
\def\red{\textsf{red}}

\theoremstyle{plain}
\newtheorem{theorem}{Theorem}
\newtheorem{conjecture}{Conjecture}
\newtheorem{lemma}[theorem]{Lemma}
\newtheorem{proposition}[theorem]{Proposition}
\newtheorem{corollary}[theorem]{Corollary}

\newtheorem*{unnumberedthm}{Main Theorem}


\theoremstyle{definition}
\newtheorem{definition}[theorem]{Definition}
\newtheorem{fact}[theorem]{Fact}
\newtheorem{remark}[theorem]{Remark}

\def\cqfd{\skip10=\parfillskip\parfillskip=0pt
\enspace\hfill\symbolecqfd\par\parfillskip=\skip10\par\medskip}
\def\symbolecqfd{\rlap{$\sqcap$}$\sqcup$}
\def\proofof#1{\rm \trivlist \item[\hskip \labelsep{\bf
Proof of~#1.}]}
\def\eopo{\cqfd\endtrivlist}

\def\notepw#1{\marginpar{\tiny\it\textcolor{orange}{#1}}}
\newcommand{\frederique}[1]{\textcolor{red}{#1}}
\newcommand{\cyril}[1]{\textcolor{blue}{#1}}
\newcommand{\pascal}[1]{\textcolor{orange}{#1}}
\let\pw\pascal
\let\fb\frederique

\title{The average-case complexity of the Word Problem for groups of matrices over $\Z$ is linear}

\author{
    % 
    Fr\'ed\'erique Bassino, \small{\url{bassino@lipn.fr}}\\
    \small{Univ. Sorbonne Paris Nord, CNRS, LIPN, UMR 7030, F-93430 Villetaneuse, France}%     
    \and
%     \\
    Cyril Nicaud, \small{\url{cyril.nicaud@univ-eiffel.fr}}\\
    \small{Univ Gustave Eiffel, CNRS, LIGM, F-77454 Marne-la-Vallée, France}%
%
    \and
%     \\
    Pascal Weil, \small{\url{pascal.weil@cnrs.fr}}\\
    \small{CNRS, ReLaX, IRL 2000, Siruseri, India}\\
    %
    \small{CNRS, Univ. Sorbonne Paris Nord, LIPN, UMR 7030, F-93430 Villetaneuse, France}%
    %     
    }
\date{\today}

\begin{document}

\maketitle

\begin{abstract}
    We show that the Word Problem in finitely generated subgroups of $\GL_d(\Z)$ can be solved in linear average-case complexity. This is done under the bit-complexity model, which accounts for the fact that large integers are handled, and under the assumption that the input words are chosen uniformly at random among the words of a given length. Our result generalizes to matrices in $\GL_d(R)$, where $R$ is a subring of $\C$, of finite rank over $\Z$.
\end{abstract}


%%%%%%%%%%%%%%%%%%%%%%%%%%%%%%%%%%%%%%
\section{Introduction}


Let $G$ be a group and let $\Sigma$ be a finite non-empty subset of $G$. The Word Problem for $G$ relative to $\Sigma$ is the following: given a finite sequence (a word) $w$ of elements of $\Sigma$ and their inverses, decide whether the value of $w$ in $G$ is trivial. This problem was introduced by Dehn in the early 20th century \cite{1911:Dehn}, and is considered one of the fundamental problems in algorithmic and combinatorial group theory. It is known that the Word Problem is not decidable in general (that is: there exist finitely generated, and even finitely presented groups with undecidable Word Problem), see Novikov \cite{1955:Novikov}
and Boone \cite{1959:Boone}. However, it is known to be decidable in many important classes of groups, for instance in automatic groups (see Epstein \emph{et al.} \cite{1992:Epsteinetal}),
which includes finite, free, hyperbolic or braid groups, in finitely presented residually finite groups (Simmons \cite{1973:Simmons}), in 1-relator groups (see Magnus \emph{et al.} \cite[Theorem 4.14]{1966:MagnusKS} and also Lyndon and Schupp \cite{1977:LyndonSchupp}), etc.

Here we consider this problem in the very natural context of subgroups of $\GL_d(\Z)$, the group of invertible matrices with integer coefficients. We set the following notation: $\Sigma$ is a finite non-empty subset of $\GL_d(\Z)$, $\alph = \Sigma \cup \Sigma^{-1}$, $H$ is the subgroup of $\GL_d(\Z)$ generated by $\Sigma$ and $\mat\colon \alph^* \to \GL_d(\Z)$ is the natural morphism, which maps the element $x \in \alph$ to the corresponding matrix in $\GL_d(\Z)$. That is, if $w$ is a word on alphabet $\alph$, then $\mat(w)$ is the value of $w$ in $\GL_d(\Z)$. The Word Problem in $H$ (relative to $\Sigma$), written $\WP_\Sigma$, is obviously decidable: $\mat(w)$ can be computed and compared to the identity matrix $\Id$.

In the following, we consider $d$ and $\Sigma$ to be fixed. The coefficients of the matrices in $H$ can be very large (the coefficients of $\mat(w)$ may grow exponentially in the length $|w|$ of $w$), and we therefore evaluate the complexity of algorithms in the so-called bit-complexity model: integers are identified with their binary expansion and arithmetic operations require more than constant time.  In particular, we use the recent and deep result of Harvey and van der Hoeven \cite{harvey2021integer} which states that multiplying two integers $p$ and $q$ can be done in time $\O(L\log L)$ where $L = \max(\log p, \log q)$ is (roughly) the maximal length of the binary expansions of $p$ and $q$.

The naive algorithm to solve the Word Problem in $H$ is the following: given $w = a_1\cdots a_n$, with each $a_i \in \alph$, let $w_0 = \Id$ (the identity matrix) and, for each $1\le i\le n$, $w_i = w_{i-1}a_i$. Then $\mat(w) = w_n$ and one can decide by inspection whether $\mat(w) = \Id$. This algorithm has a worst-case complexity in $\O(n^2\log n)$. A direct application of the classical divide-and-conquer strategy lowers this complexity to $\O(n\log^2 n)$, as noted already by Olshanskii and Shpilrain \cite[Proposition 2]{2025:OlshShp} (see \cref{prop: complexity DC} below).

We note that this divide-and-conquer technique yields an $\O(n)$ worst-case complexity when the matrices in $\Sigma$ are upper-triangular (\cref{prop: upper triangular}), a sharpening of \cite[Theorem 2]{2025:OlshShp}. As noted in that paper, this implies a linear worst-case complexity for the Word Problem in finitely generated nilpotent groups.

Our main result (see \cref{thm: main theorem} below for a more precise statement) is the following.

\begin{unnumberedthm}
    %
    $\WP_\Sigma$ has linear time average-case complexity in the bitcost model of computation when inputs follow the uniform distribution on length $n$ words over alphabet $\alph$.
    %
\end{unnumberedthm}

This result was already known in the case of polycyclic groups, which are representable as subgroups of $\GL_d(\Z)$ (see Wehrfritz \cite{1980:Wehrfritz}). Indeed, Olshanskii and Shpilrain established a linear average-case complexity of the Word Problem for these groups \cite[Theorem 3]{2025:OlshShp}, and in fact, for subgroups $H \le \GL_d(\Z)$ with a non-trivial virtually abelian factor, such as the virtually solvable linear groups \cite[Remark 4]{2025:OlshShp}, which goes beyond polycyclic groups.

In contrast, our result holds for all finitely generated subgroups of $\GL_d(\Z)$, polycyclic or not, and its proof is seemingly very different from that in \cite{2025:OlshShp}. The two proofs have however an interesting common point, see \cref{rem: compare quotienting proofs} below.

Several ideas come into play in the proof of our main result. The first is to use modulo computations: if $q$ is an integer, one first solves the Word Problem in the subgroup $H_q$ of $\GL_d(\Z/q\Z)$ generated by $\Sigma$. That can be done by using the standard (divide-and-conquer) algorithm in $\GL_d(\Z/q\Z)$: that is, by computing $\mat(w)_q$ (the matrix $\mat(w)$ modulo $q$) and verifying whether it is equal to $\Id$. Here we benefit from the fact that, at each step of the computation, the entries of matrices are in the interval $[0, q-1]$, thus lowering the worst-case complexity. If $q$ is constant, this is faster than computing $\mat(w)$ in $\GL_d(\Z)$ as the length of $w$ tends to infinity; and if $\mat(w)_q$ is not the identity, then neither is $\mat(w)$. Thus, if $\mat(w)_q \ne \Id$ with high probability, we first compute $\mat(w)_q$ and, if it is the identity, we compute $\mat(w)$ using the standard $\O(n\log^2 n)$ algorithm.

This however is not sufficient to get our result.
%
Indeed, as we will see, for any \emph{fixed} value of $q$, the probability that $\mat(w)_q = h$ (for any $h \in \GL_d(\Z/q\Z)$) may tend to a positive value, namely $\frac\alpha{|H_q|}$, where $\alpha = 1$ or $\alpha = 2$ (depending on $\Sigma$ and $q$). Thus, with probability tending to $\frac\alpha{|H_q|}$, we need to call the divide-and-conquer algorithm, yielding an $\O(n\log^2n)$ average-case complexity.

The next idea is to choose the modulus $q$ as a function $q(n)$ of the length $n$ of the input word $w$. The advantage of this algorithm is that, for each input word $w$, we compute $\mat(w)$ modulo a single integer $q(n)$. The challenge is to identify an appropriate function $q(n)$.

More precisely, we see the computation of $\mat(w)_{q(n)}$ in terms of trajectories in a Markov chain $\tM$, which is a technical variant of the natural Markov chain based on the Cayley graph of the subgroup $H_{q(n)}$ of $\GL_d(\Z/q(n)\Z)$ generated by $\Sigma$. We want $q(n)$ to be small enough so that $\mat(w)_{q(n)}$ is computed rapidly, and we want the Markov chain $\tM$ to have low mixing time (that is, we want it to converge rapidly towards its stationary distribution) and a large dispersion (so that $H_{q(n)}$ should have large cardinality and the probability that $\mat(w)_{q(n)} = \Id$ be small). These are seemingly contradictory requirements, and the technical work of the proof consists in identifying a function $q(n)$ with these properties.

Concretely, we choose $q(n)$ to be an increasingly long product of distinct primes, so that we test $\mat(w)$ modulo all these primes in a single and quick computation (namely, the modulo $q(n)$ computation), yet $q(n)$ grows sufficiently slowly. See \cref{def: def of q} for more details.


\begin{remark}\label{rem: compare quotienting proofs}
    The proof of \cite[Theorem 3]{2025:OlshShp} also uses computation in a quotient of $H$, though not a $\bmod$ $q$ quotient: the projection onto an abelian factor, whose existence is postulated (and which should be known for any implementation of their algorithm).
\end{remark}

It is important to note that our algorithm makes no assumption on $\Sigma$ (except the fact that $\Sigma \cap \Sigma\inv = \emptyset$) and on the properties of the subgroup $H$ it generates. The matrices in $\Sigma$ may be triangular or not; the subgroup $H$ may be finite (and it is well known that the Word Problem is simpler in this case) or infinite; it may be nilpotent, polycyclic or virtually solvable; it may have exponential or polynomial growth (as a consequence of the Tits alternative): we do not need to identify in which situation we are, and we always use the same algorithm.

The paper is organized as follows. In \cref{sec: preliminaries}, we lay out the notation and definitions of the Word Problem and its modulo variant and we give a very quick discussion of models of computation and of the definitions of worst-case and average-case complexity. In \cref{sec: computing with integers}, we remind the readers of the precise complexity of computing with integers, and in \cref{sec: markov chains}, we recall the fundamental definitions and results on Markov chains and their convergence to a stationary distribution, inasmuch as they will be needed in this paper.

In \cref{sec: algorithms}, we describe and analyze the standard divide-and-conquer algorithm to compute $\mat(w)$ (\cref{sec: first algorithms}), and we briefly discuss its application in the case of triangular matrices (\cref{sec: triangular}). Our Algorithm $\S$, which solves the Word Problem in linear average-case complexity, is given in \cref{sec: better algorithm}. This includes defining the function $q(n)$ mentioned above. We also reduce the proof of our Main Theorem to a technical statement (\cref{thm: technical theorem}), namely the fact that, if $H$ is infinite, then the probability that a word % (resp. reduced word) 
$w$ of length $n$ satisfies $\mat(w)_{q(n)} = \Id$ is $\O(\log^{-2}n)$. Finally, the proof of \cref{thm: technical theorem} is given in \cref{sec: non reduced words}.

The concluding remarks in \cref{sec: concluding} include a discussion of an alternative probabilistic algorithm to solve the Word Problem, which does not result in an improved asymptotic performance. We also explain there how our result extends to matrices over a subring $R$ of the complex field $\C$, provided that the additive group $(R,+)$ is finitely generated.


%%%%%%%%%%%%%%%%%%%%%%%%%%%%%%%%%%%%%%
\section{Preliminaries and notation}\label{sec: preliminaries}

If $A$ is a matrix in $\GL_d(\Z)$, we let $\|A\|_\infty = \max\{|A_{i,j}| \mid 1\le i,j\le d\}$. If $m\ge 2$, $A_m$ denotes the projection of $A$ modulo $m$, a matrix with entries in $\Z/m\Z$.

If $X$ is a set of matrices in $\GL_d(\Z)$, the subgroup they generate is written $\langle X\rangle$. We also denote by $X_m$ the set $\{A_m \mid A\in X\}$.
%
It is interesting to note the following elementary fact.

\begin{fact}\label{rk: mod m, gl(Z) is a group}
    Let $m\ge 2$. If $m$ is not a prime, then $\Z/m\Z$ is not a field. However, the projection mod $m$ of the group $\GL_d(\Z)$ (or of any of its subgroups) is again a group.
\end{fact}



%%%%%%%%%%%%%%%%%%%%%%%%%%%%%%%%%%%%%%
\subsection{The Word Problem and related algorithmic problems}\label{sec: WP and related pbms}

Let $\Sigma$ be a fixed, finite, non-empty set of matrices in $\GL_d(\Z)$, let $\alph = \Sigma \cup \Sigma^{-1}$ and let $\alph^*$ be the set of all words on alphabet $\alph$ (i.e., finite sequences of elements of $\alph$). We denote by $\mat\colon \alph^* \to \GL_d(\Z)$ the natural (monoid) morphism, which maps each element of $\Sigma$ to itself.

The \emph{Word Problem} $\WP_\Sigma$, is the following: given a word $w \in \alph$, decide whether $\mat(w)$ is the identity matrix $\Id$.

We also consider in this paper the closely related \emph{Exact Computation Problem} $\EC_\Sigma$: on input a word $w\in \alph^*$, $\EC_\Sigma$ computes the product $\mat(w)$ of that sequence in $\GL_d(\Z)$.
%
We will also discuss the same problems \emph{modulo $m$}, where $m\ge 2$ is an integer. More precisely, Problems $\WP_{\Sigma,m}$ (resp. $\EC_{\Sigma,m}$) takes a word $w\in  \alph$ as input, and decides whether $\mat(w)_m$ is the identity matrix (resp. computes $\mat(w)_m$).

\begin{remark}\label{rk: from EC to WP}
    %
    If an algorithm $\calS$ solves $\EC_\Sigma$ (resp. $\EC_{\Sigma,m}$) on input $w$ --- that is, if we have computed $\mat(w)$ (resp. $\mat(w)_m$), --- then a minor tweak solves $\WP_\Sigma$ (resp. $\WP_{\Sigma,m}$): it suffices to examine the $d^2$ entries of $\mat(w)$ (resp. $\mat(w)_m$), which is done in constant time, under the reasonable assumption that the length of the binary representation of an integer can be compared to 1 in constant time.
    %
\end{remark}

\paragraph{Convention}
%
Throughout the paper, the integer $d$ and the set $\Sigma$ are fixed. We let $k = |\Sigma|$. We also assume, and this is no loss of generality, that $\Sigma$ does not contain a matrix $A$ and its inverse. In particular, it does not contain the identity matrix $\Id$, and $|\alph| = 2|\Sigma|$.


%%%%%%%%%%%%%%%%%%%%%%%%%%%%%%%%%%%%%%
\subsection{About algorithms and complexity}

To evaluate the complexity of an algorithm $\calS$, we consider the function $S(w)$, where $w$ is an input word, which measures the time (number of elementary operations) needed to run Algorithm $\calS$ on input $w$. Observe that, even though $\Sigma$ --- and thus the coefficients of the matrices in $\alph$ --- is fixed in our setting, the integers computed by the algorithms can be huge, and one cannot assume that arithmetic operations are performed in constant time. To account for this, we consider the \emph{bit-cost} model of computation, where an integer $n$ is encoded using roughly $\log n$ bits (see \cref{sec: computing with integers}).

The worst-case complexity of $\calS$ is the function on integers given by
%
$$S_{\textsf{wc}}(n) = \max\{S(w) \mid |w| = n\},$$
%
and its average complexity $S_{\textsf{ac}}(n)$ is the average value of $S(w)$, when $w$ runs uniformly over inputs of length $n$. As is traditional, these functions are considered up to asymptotic equivalence, when $n$ tends to infinity.

As mentioned earlier, our average-complexity results assume that input words are taken uniformly at random among words on $\alph$ of length $n$.


%%%%%%%%%%%%%%%%%%%%%%%%
\subsection{Computing with integers}\label{sec: computing with integers}
%
The \emph{length} (or \emph{bit-size}) $\ell(n)$ of a integer $n$ is the length of its binary expansion, namely $\ell(n) = \lceil\log (|n|+1)\rceil + 1$ (all logarithms are in base 2), where the additional bit is used to encode the sign. As a result, the integers of length at most $\ell + 1$ have absolute value less than $2^\ell$. We freely use the following facts, all of them elementary --- with the exception of \cref{prop: basic facts}~\eqref{fact: HvdH}, which is a deep result due to Harvey and van der Hoeven \cite{harvey2021integer}.

\begin{proposition}\label{prop: basic facts}
%
Let $n,n'$ be integers and let $L$ such that $\ell(n), \ell(n') \le L$.
%
\begin{enumerate}[(i)]
%
\item $\ell(nn') \le \ell(n) + \ell(n') - 1$. 
%
\item\label{fact: HvdH} The product $nn'$ is computed in time $\O(L\log L)$.

\item The product $nn'$ is also computed (by the primary school multiplication algorithm) in time $\O(\ell(n)\ell(n'))$, which is interesting if $\ell(n)$ is small with respect to $\ell(n')$.

\item The sum of $d$ integers of length at most $L$, has length at most $1 + \log d + L$ and is computed in time $\O(d\log d + dL)$.

\item\label{fact: product by general matrix} If $A, A'$ are $d \times d$ matrices with entries of length at most $L$, the entries of the product $AA'$ have length at most $1 + \log d + 2L$, and each is computed in time $\O(d\log d + dL\log L)$.

\item\label{fact: product by fixed matrix} If the entries of $A$ (resp. $A'$) are of length at most $L$ (resp. $L'$), and if $L'$ is much smaller than $L$, the entries of the product $AA'$ have length at most $1 + \log d + L + L'$, and each is computed in time $\O(d\log d + dLL')$.
%
\end{enumerate}
%
\end{proposition}

\begin{remark} \label{rk:coef growth}
\cref{prop: basic facts}~\eqref{fact: product by fixed matrix} shows that, if $w$ is a length $n$ word on $\alph$, the entries of $\mat(w)$ have length at most $(1+L+\log d)\,n$, where $L$ is the maximum length of the coefficients of the matrices in $\alph$: the length of the entries of $\mat(w)$ grows at most linearly in $n$, and their value in $\Z$ at most exponentially.
\end{remark}

In the sequel, we will also compute in $\Z/m\Z$ for some integer $m\ge 2$. We record the following result on the complexity of computing in this ring, where every element is represented by a non-negative integer at most equal to $m$. It follows directly from \cite[Theorem 9.8 and Corollary 9.9]{vzGathenGerhard}, together with Proposition~\ref{prop: basic facts}~\eqref{fact: HvdH} due to \cite{harvey2021integer}.

\begin{corollary}\label{cor: arithmetic mod p}
    Let $m\ge 2$. Every arithmetic operation in $\Z/m\Z$ is performed in $\O(\log m\, \log\log m)$.
\end{corollary}


%%%%%%%%%%%%%%%%%%%%%%%%
\subsection{Basic results on probability distributions and Markov chains}\label{sec: markov chains}

For a general discussion of probability distributions and Markov chains, we refer readers to \cite{LevinPeres}. Here we fix some notation and state a few standard results on Markov chains, that will be used in the sequel.

If $X$ is a set, a \emph{probability distribution} (or \emph{probability vector}) on $X$ is a vector $(\mu(x))_{x\in X}$, where each $\mu(x)$ lies in the closed interval $[0,1]$ and $\sum_x\mu(x) = 1$. The \emph{corresponding probability function} is given, for every subset $Y$ of $X$, by $\mu(Y) = \sum_{y\in Y}\mu(y)$. The \emph{uniform distribution} on $X$ is the vector all of whose entries are $\frac1{|X|}$.

If $\mu$ and $\pi$ are two probability distributions on the same finite set $X$, their \emph{total variation distance} is 
%
\begin{equation}\label{eq:var}
\|\pi - \mu\|_{\text{Var}} = \sup_{A\subseteq X}\,|\pi(A) - \mu(A)| = \frac12\sum_{x\in X}|\pi(x)-\mu(x)|.
\end{equation}
%
The second equality is proven in \cite[Proposition 4.2]{LevinPeres}.

In this paper, a (finite) \emph{Markov chain} $\tM$ consists in a directed graph with (finite) vertex set $S$ (vertices are also called \emph{states}), edge set a subset of $S\times S$, and, for each edge (called a \emph{transition}) from state $q$ to state $q'$, of a value $M_{q,q'} \in [0,1]$, in such a way that, for each $q\in S$, $(M_{q,q'})_{q'\in S}$ is a probability vector. The matrix $M = (M_{q,q'})_{(q,q') \in S \times S}$ is called the \emph{transition matrix} of $\tM$.

A path $T = (q_0,\dots,q_n)$ in the underlying graph of $\tM$ is called a \emph{trajectory}. A probability is assigned to $T$ by $\tM$, namely the product $M_{q_0,q_1}\,M_{q_1,q_2}\dots\,M_{q_{n-1},q_n}$. We note that the $(q,q')$-entry of the $n$-th power of the transition matrix $M$ is the sum of the probabilities of the length $n$ trajectories from $q$ to $q'$.

The Markov chain $\tM$ is \emph{symmetric} if its transition matrix is symmetric, that is, if $M_{q,q'} = M_{q',q}$ for all states $q,q'$. The chain $\tM$ is \emph{irreducible} if its underlying graph is strongly connected, that is: for every $q, q'\in S$, the $(q,q')$-entry of some positive power of $M$ is non-zero. The chain $\tM$ is said to be \emph{aperiodic} if, for every $q \in S$ and for all $n$ large enough, there exists a length $n$ trajectory from $q$ to $q$. Finally, the chain $\tM$ is \emph{primitive} if it is both irreducible and aperiodic.

%It is directly verified that the eigenvalues of the transition matrix $M$ have modulus less than or equal to 1. 
A probability distribution $D$ on $X$ is called \emph{stationary} if $D\,M = D$, that is, if $D$ is a left eigenvector for the eigenvalue 1.

A Markov chain $\tM$ with a stationary distribution $\pi$ is said to be \emph{reversible with respect to $\pi$} (or just \emph{reversible} if the stationary distribution is unique) if $\pi(q)\,M(q,q') = \pi(q')\,M(q',q)$ for all states $q,q'$.

The following is a classical result on Markov chains.

\begin{theorem}\label{thm: basics on Markov}
    Let $\tM$ be a primitive Markov chain. Then 1 is an eigenvalue of $M$ and the other eigenvalues have modulus less than 1. The eigenspace corresponding to eigenvalue 1 has dimension 1 and $\tM$ has a unique stationary distribution $\pi$, which satisfies the following: if $\mu$ is any probability distribution on the state set of $\tM$, then $\lim_n \mu\,M^n = \pi$. 

    If $\tM$ is primitive and symmetric, then the uniform distribution is its unique stationary distribution, all the eigenvalues are real, and $\tM$ is reversible.
\end{theorem}


%%%%%%%%%%%%%%%%%%%%
\section{Algorithms for $\WP_\Sigma$ and $\EC_\Sigma$}\label{sec: algorithms}

With the aim of studying the average-case complexity of $\WP_\Sigma$, we consider several algorithms solving this problem. We start with standard algorithms, including one with $\O(n\log^2n)$ worst-case complexity (\cref{sec: first algorithms}). In \cref{sec: better algorithm}, we introduce a better algorithm with the announced linear average-case complexity for  the uniform distribution on words of length $n$.
We then state our main theorem, \cref{thm: main theorem}, and reduce its proof to a technical statement (\cref{thm: technical theorem}).
The proof of that statement is given in \cref{sec: non reduced words}.

%%%%%%%%%%%%%%%%%%%%
\subsection{First algorithms}\label{sec: first algorithms}
%

The naive algorithm to compute $\mat(w)$ consists in reading the word $w$ from left to right, one letter at a time, and performing the corresponding $n-1$ matrix multiplications --- where $n = |w|$. The right factor in each of these operations is a matrix in $\alph$, a constant set which is independent of $n$. A direct application of \cref{prop: basic facts}~\eqref{fact: product by fixed matrix} and \cref{rk:coef growth} then shows that the worst-case bit complexity of this algorithm is $\O(\ell^2n^2)$, where $\ell$ is an upper bound of the bit-size of the coefficients of the elements of $\alph$. Since $\ell$ is fixed in our settings, this naive approach runs in $\O(n^2)$ time.

This quadratic upper bound can be significantly improved using a \emph{divide and conquer} strategy.

\medskip

\begin{algorithm}[H]
\caption{\bf Algorithm $\DC_\Sigma$}
\DontPrintSemicolon
\SetKwInOut{Input}{Input}
\SetKwInOut{Output}{Output}

\Input{a sequence $w$ of $n$ elements of $\alph$}
\Output{$\mat(w)$}

\lIf{$n=0$ (resp. $n=1$)}{\Return $\Id$ (resp. $\mat(w)$)}
$w_1\gets$ prefix of $w$ of length $\lfloor n/2\rfloor$\;
$w_2\gets$ suffix of $w$ of length $\lceil n/2\rceil$\;
\Return $\DC_\Sigma(w_1)\times\DC_\Sigma(w_2)$
\end{algorithm}

\medskip

Before we analyze the complexity of this simple algorithm, let us remind the reader of an instance of the celebrated \emph{Master Theorem} (see, \textit{e.g.}, \cite[Theorem 4.1]{cormen2022introduction}), which we will use several times.

\begin{proposition}\label{prop: master theorem}
    %
    Let $C(n)$ and $f(n)$ be positive-valued non-decreasing functions on $\N$ satisfying the equation $C(n) = C(\lfloor \frac n2\rfloor) + C(\lceil \frac n2\rceil) + f(n)$ for $n\geq 2$.
    %
    \begin{itemize}
        %
        \item If $f(n) = \O(n^h)$ for some $0 \leq h < 1$, then $C(n) = \O(n)$.
        %
        \item If $f(n) = \O(n\,\log^hn)$ for some $h \geq 0$, then $C(n) = \O(n\,\log^{h+1}n)$.
        %
    \end{itemize}
    %
\end{proposition}    

\begin{proposition}\label{prop: complexity DC}
%
Algorithm $\DC_\Sigma$ solves Problem $\EC_\Sigma$, with worst-case complexity $\O(n\log^2n)$. In addition, Problem $\WP_\Sigma$ can be solved with the same worst-case complexity. 

Moreover, if $H = \langle\Sigma\rangle$ is finite, then the complexity of $\DC_\Sigma$ is linear.
%
\end{proposition}

\begin{proof}
%
Since $w = w_1w_2$, we have $\mat(w) = \mat(w_1)\,\mat(w_2)$, so Algorithm $\DC_\Sigma$ solves Problem $\EC_\Sigma$. 

\cref{prop: basic facts}~\eqref{fact: HvdH} (crucially using~\cite{harvey2021integer}) and Remark~\ref{rk:coef growth} show that the complexity $C(n)$ of this algorithm satisfies the equation

\begin{equation}
    C(n) = C\left(\lfloor n/2\rfloor\right) + C\left(\lceil n/2\rceil\right) + \O(n\log n) \qquad\text{for } n\geq 2.
\end{equation}
%
\cref{prop: master theorem} then yields the fact that $C(n)$ is $\O(n\log^2n)$. The statement on Problem $\WP_\Sigma$ follows from \cref{rk: from EC to WP}.

If $H$ is finite, then the coefficients of the matrices in $H$ have bounded length and the complexity $C(n)$ now satisfies the equation
\begin{equation}
    C(n) = C\left(\lfloor n/2\rfloor\right) + C\left(\lceil n/2\rceil\right) + \O(1) \qquad\text{for } n\geq 2.
\end{equation}
%
\cref{prop: master theorem} then yields the fact that $C(n)$ is $\O(n)$.
%
\end{proof}

The same algorithm can be run on matrices in $\Z/m\Z$, where $m \ge 2$ is any integer. Let $\DC_{\Sigma,m}$ be the algorithm with the same steps as Algorithm $\DC_\Sigma$, where all arithmetic operations are performed in $\Z/m\Z$ instead of $\Z$. It is immediate that Algorithm $\DC_{\Sigma,m}$ solves Problem $\EC_{\Sigma,m}$, and hence also Problem $\WP_{\Sigma,m}$ in linear time (when $m$ is fixed).

\begin{remark}\label{rk: worst case in other rings}
    In fact, the same algorithm runs on matrices over any computable ring. Over $\Q$, it also yields an $\O(n\log^2n)$ worst-case complexity since the addition and multiplication of rationals takes asymptotically the same time as the addition and multiplication of integers.
\end{remark}



%%%%%%%%%%%%%%%%%%%%
\subsection{The special case of triangular matrices}\label{sec: triangular}

We note that if $A$ and $A'$ are upper triangular matrices and if $1\le i \le j \le d$, then the $(i,j)$-entry of $AA'$ is $\sum_{h=i}^jA_{i,h}\,A'_{h,j}$. 

Now suppose that the matrices in $\Sigma$ are upper-triangular, and hence so are their inverses. It is directly verified that, if $w \in \alph^*$ has length $n$ and $1\le i \le j \le d$, then the $(i,j)$-entry of $\mat(w)$ is bounded above by a polynomial in the variable $n$, with degree $j-i$. In particular, the length of the entries of $\mat(w)$ is logarithmic.

Then \cref{prop: basic facts} shows that if $w$ and $w'$ are words of length at most $n$, then the product $\mat(w)\mat(w')$ is computed in polylogarithmic time. Applying the Master Theorem (\cref{prop: master theorem}) directly yields a linear worst-case complexity for $\DC_\Sigma$ and $\WP_\Sigma$. As noted by Olshanskii's and Shpilrain \cite{2025:OlshShp}, this result can be combined with the well known fact \cite{1938:Hirsch} that any finitely generated nilpotent group can be embedded in the direct product of a finite group and a group of upper-triangular matrices in some $\GL_d(\Z)$: this yields the following statement, which improves on \cite[Theorems 1 and 2]{2025:OlshShp}.

\begin{proposition}\label{prop: upper triangular}
    If $\Sigma$ consists only of upper (resp. lower) triangular matrices, then the worst-case complexity of $\WP_\Sigma$ is $\O(n)$.

    If $G$ is a finitely generated nilpotent group, then the Word Problem in $G$ can be solved in linear worst-case complexity.
\end{proposition}


%%%%%%%%%%%%%%%%%%%%
\subsection{A linear average-case algorithm for the Word Problem}\label{sec: better algorithm}

The key idea to get an algorithm solving $\WP_\Sigma$ with a better average-case complexity, is to compute $\mat(w)_{q(n)}$, where $n$ is the length of $w$ and $q(n)$ is a function such that 
\begin{enumerate}[(i)]
    \item $\mat(w)_{q(n)}$ is unlikely to be the identity, and
    \item $\mat(w)_{q(n)}$ can be computed in linear time. 
\end{enumerate}
%
In the rest of the paper, we use the following function $q$.

\begin{definition}\label{def: def of q}
    Let $q\colon \N\to\N$ be the function given by $q(0) = q(1) = 1$ and, for all $n\ge 2$,
\[
q(n) = \prod_{\substack{p \enspace \leq \enspace \log^5 n\\p\text{ prime}}} p
\]
\end{definition}

Algorithm $\S$ is the following. 

\medskip

\begin{algorithm}[H]
\caption{\bf Algorithm $\S$}
\DontPrintSemicolon
\SetKwInOut{Input}{Input}
\SetKwInOut{Output}{Output}

\Input{a sequence $w$ of $n$ elements of $\alph$}
\Output{\True\ if $\mat(w) = \Id$, and \False\ otherwise}

Compute $q(n)$\;
\eIf{$\DC_{\Sigma,q(n)}(w)\neq\Id$}{\Return \False}
{\eIf{$\DC_{\Sigma}(w) \neq \Id$}{\Return \False}{\Return \True}}
%
\end{algorithm}

\medskip

We can now give a precise version of our main theorem, stated in the introduction.

\begin{theorem}\label{thm: main theorem}
    %
    Algorithm $\S$ solves Problem $\WP_\Sigma$ with linear time average-case complexity, when inputs are uniform random words.
    %
\end{theorem}

The reader should note that Algorithm $\S$ makes no assumption on the algebraic or combinatorial properties of $\Sigma$ or the subgroup $H = \langle\Sigma\rangle$. The same algorithm is run, with linear average-case complexity, whether $\Sigma$ consists of triangular matrices or not, and whether $H$ is finite or infinite. The latter property is decidable (Jacob, \cite{1978:Jacob}) in polynomial time (Babai, Beals and Rockmore \cite{1993:BabaiBealsRockmore}, see also Detinko and Flannery \cite{2009:DetinkoFlannery}). Similarly, the same algorithm is run, with the same average-case complexity whether $H$ has polynomial or exponential growth, or whether it is nilpotent (see \cref{sec: triangular}), polycyclic or virtually solvable. In the latter two situations, Olshanskii and Shpilrain recently proved a linear average-case complexity of the Word Problem \cite[Theorem 3 and Remark 4]{2025:OlshShp}, using the properties of these subgroups, namely computing in a quotient satisfying a specific algebraic condition. 

Towards the proof of \cref{thm: main theorem}, we record the following technical statements.

\begin{proposition}\label{prop: t polylogarithmic}
    The function $q$ can be computed in polylogarithmic time, and $q(n)$ has polylogarithmic length.
\end{proposition}

\begin{proof}
    %
A rough upper bound for $q(n)$ is $(\log^5 n)^{\log^5 n}$, so that $\log q(n)$ (and hence the length of $q(n)$) is bounded above by $5\log^6n$, a polylogarithm.

It follows from algorithms proposed by Mairson \cite{mairson1977some} and Pritchard \cite{1987:Pritchard} that one can list all prime numbers less than or equal to $N$ in (bit-complexity) $\O(N\, \log N\, \log\log N))$. In particular, listing the prime numbers at most equal to $\log^5n$ is done in $\O(\log^5 n\, \log\log n \,  \log\log\log n)$.


Computing the product of two numbers at most equal to $q(n)$ (and hence with length at most $5 \log^6n$) is done in time $\O(\log^6n\, \log\log n)$ by \cref{prop: basic facts}~\eqref{fact: HvdH}. It follows that $q(n)$, the product of at most $\log^5n$ numbers less than or equal to $q(n)$, is computed in time $\O(\log^{11}n\log\log n)$: again a polylogarithm. (Using a divide and conquer method yields a polylogarithm with lesser degree.)
    %
\end{proof}

\begin{lemma}\label{lemma: ac of DCm}
%
Let $P_n$ be the probability that a word $w$ of length $n$ satisfies $\mat(w)_{q(n)} = \Id$. The aver\-age-case complexity of $\S$ (when inputs are uniform random words of length $n$) is $\O(n + P_n\, n\, \log^2n)$.

Moreover if $H = \langle\Sigma\rangle$ is finite, then the average-case complexity of $\S$ is $\O(n)$.
%
\end{lemma}

\begin{proof}
    %
    Once $q(n)$ is computed (in polylogarithmic time, by \cref{prop: t polylogarithmic}), the complexity $C(n)$ of the second step of Algorithm $\S$, that is, of running $\DC_{\Sigma,q(n)}$ on a word $w$ of length $n$, satisfies
    %
    $$C(n) = C\left(\lfloor n/2\rfloor\right) + C\left(\lceil n/2\rceil\right) + \O(\log q(n)\,\log\log q(n))\qquad \text{for } n\geq 2$$
    %
    since every arithmetic operation in $\ZZ/q(n)\ZZ$ can be performed in time $\O(\log q(n)\, \log\log q(n))$ (see \cref{cor: arithmetic mod p}). By \cref{prop: t polylogarithmic} again, we have $\log q(n)\,\log\log q(n) =o(n)$, and 
    \cref{prop: master theorem} yields the fact that $C(n)$ is linear.
    
    Finally, by \cref{prop: complexity DC}, the worst-case complexity of Algorithm $\DC_{\Sigma}$ on inputs of length $n$ is in $\O(n\log^2 n)$ if $H$ is infinite and is linear if $H$ is finite.
    %
    The announced result follows.
%
\end{proof}


In \cref{sec: non reduced words}, we prove the following statement.

\begin{theorem}\label{thm: technical theorem}
    %
    If $H = \langle\Sigma\rangle$ is infinite, then $\mat(w)_{q(n)} = \Id$ with probability $\O(\log^{-2} n)$ (when $w$ is chosen uniformly at random among length $n$ words).
    %
\end{theorem}

The proof of \cref{thm: main theorem} follows directly.

\begin{proof}[Proof of \cref{thm: main theorem}]
    %
    It is immediate that Algorithm $\S$ solves $\WP_\Sigma$. Note that the subgroup $H$ is fixed in our setting (that is: it is not part of the input). If $H$ is finite, the result was established in \cref{lemma: ac of DCm}.
    
    If $H$ is infinite, then \cref{thm: technical theorem} and \cref{lemma: ac of DCm} show that $\S$ runs with linear average-case complexity.
    %
\end{proof}

We are now left with proving \cref{thm: technical theorem}: this is done in \cref{sec: non reduced words}.

%%%%%%%%%%%%%%%%%%%%%%%%%%%%%%%%%%%
\section{Proof of \cref{thm: technical theorem}}\label{sec: non reduced words}

Recall that $H$ denotes the subgroup of $\GL_d(\Z)$ generated by $\Sigma$, and that $H_m$ denotes the mod $m$ projection of $H$.

Throughout this section, we let $m$ be an integer greater than $2\,\max_{A\in \alph}\|A\|_\infty$. We note that $q(n)$ satisfies this condition for all $n$ large enough. It follows from this assumption that the matrices $A_m$ ($A\in \alph$) are pairwise distinct and, in particular, none is the identity modulo $m$.


%%%%%%%%%%%%%%%%%%%%%%%%%%%%%%%%%%%%%%%%%
\subsection{Uniform random words as trajectories in a Markov chain}

The matrices $\mat(w)_m$, when $w$ is a random word of length $n$, are naturally produced by the length $n$ trajectories in the Markov chain $\U_m$ defined below.
Recall that $|\Sigma| = k$.
%
\begin{enumerate}[(i)]
%
\item The state set of $\U_m$ is the subgroup $H_m$.
%
\item\label{item: edge labels} There is an edge $M\xrightarrow{\frac{1}{2k}} M'$ if and only if there exists a matrix $A\in\alph$ such that $M\, A = M'$.
%
\item The initial vector assigns probability 1 to $\Id$ and probability 0 to all the other states.
%
\end{enumerate}
%
Let $P_m$ be the transition matrix of $\U_m$. We formulate the following elementary observations.
%
\begin{itemize}
    %
    \item In Item~\eqref{item: edge labels}, the matrix $A \in \alph$ is uniquely determined (if it exists) by the origin $M$ and the end $M'$ of the edge: we will denote this edge by $M\xrightarrow{A:\frac{1}{2k}} M'$ when needed. We call $A$ \emph{the matrix label} of that edge.
    %
    \item The underlying graph of $\U_m$ is strongly connected, that is, the Markov chain $\U_m$ is irreducible.
    %
    \item The Markov chain $\U_m$ is symmetric, since $\alph$ contains both the matrices in $\Sigma$ and their inverses.
    
    \item In the underlying graph of $\U_m$, replacing each edge label of the form $A:\frac{1}{2k}$ by its matrix label $A$, yields the Cayley graph of the subgroup $H_m$. By the assumption made on $m$, no edge of $\U_m$ is a loop. 
    %
    \item If $T_n$ is a length $n$ trajectory in $\U_m$ and $\lab(T_n)$ denotes the word obtained from $T_n$ by reading the sequence of matrix labels of the edges traversed by $T_n$, then $w = \lab(T_n)$ is a uniform random word of $\alph^n$, and $T_n$ ends at state $\mat(w)_m$.
    %
    \item Let $\vec{\bf 1}$ be the (column) vector all of whose entries are 1. Then $P_m\, \vec{\bf 1} = \vec{\bf 1}$, so the uniform distribution vector $\frac1{|H_m|}\vec{\bf 1}$ is a right eigenvector for the eigenvalue $1$. Since $P_m$ is symmetric, it is also a left eigenvector for the eigenvalue $1$. 
    %
\end{itemize}

As observed, the Markov chain $\U_m$ is irreducible, but it may not be aperiodic. However, $\U_m$ contains cycles of length two (for instance $\Id\xrightarrow{A}A_m\xrightarrow{A\inv}\Id$ for any $A\in \Sigma$), and since its period is the gcd of the lengths of its cycle, it is equal to  $1$ or $2$. To deal with both cases at once, we consider the Markov chain $\U^2_m$, which performs two consecutive steps in $\U_m$, and whose transition matrix is $P_m^2$. More precisely, if $M, M'\in H_m$, then
\[
P^2_m(M,M') = \sum_{\substack{A,B\in\alph\\M'=MAB}} \frac1{4k^2}
\]
%
(where products are taken in $H_m$). Given our hypothesis on $m$, for each $A\in \alph$, there is at most one matrix $B\in \alph$ such that $MAB = M'$, and hence $P^2_m(M,M')\leq \frac1{2k}$.

In the particular case where $M' = M$, for each $A\in \alph$, we have $MAA\inv = M$, so $P^2_m(M,M) = \sum_{A\in\alph} \frac1{4k^2} = \frac1{2k}$.

As it contains self-loops, the Markov chain $\U^2_m$ is aperiodic. However, if $\U_m$ has period 2, then $\U^2_m$ is not strongly connected. More precisely, if $\U_m$ has period 2, then $\U_m^2$ is the disjoint union of two chains with the same number of states: one for the states at an even distance from $\Id$ in $\U_m$, and one for the states at an odd distance from $\Id$. The state set of the former is $\{M(w)_m:w\in\alph^*\text{ and }|w|\text{ even}\} = \langle \alph^2\rangle$, and it has cardinality $\frac12|H_m|$.

Let $\tU_m$ be the restriction of $\U_m$ to the set $\tH_m$ of states that are accessible from $\Id$ in $\U_m^2$. Again: if $\U_m$ is aperiodic, then $\tU_m = \U_m$ and $\tH_m = H_m$; and if $\U_m$ has period 2, then $|\tH_m|=\frac12|H_m|$. We let $\tP_m$ be the restriction of $P_m^2$ to $\tH_m$, that is,
$\tP_m(M,M')=P_m^2(M,M')$ for all $M,M'\in\tH_m$.

To summarize, we have the following statement.

\begin{proposition}\label{prop: summary tUm}
    $\tU_m$ is a symmetric and primitive Markov chain, whose set of states $\tH_m$ satisfies $|\tH_m|\geq \frac12|H_m|$. The uniform distribution on $\tH_m$ is its unique stationary distribution and $\tU_m$ is also reversible. For every $M\in \tH_m$, we have $\tP_m(M,M)=\frac1{2k}$. Moreover, for all $M,M'\in\tH_m$ such that $\tP_m(M,M')>0$, we have $\frac1{4k^2} \leq \tP_m(M,M') \leq \frac1{2k}$.
\end{proposition}


%%%%%%%%%%%%
\subsection{Probability of being the identity in $\tU_m$}

By \cref{prop: summary tUm}, the distribution of the states reached after a random length $n$ trajectory, starting at any state in $\tH_m$, converges to the uniform distribution $\pi$.
%
We now need to evaluate the rate of this convergence.
%We first establish a lemma on length $n$ trajectories in the Markov chain $\tU_m$.

\begin{lemma}\label{lm:convergence rate}
Let $m$ be an integer greater than $2\,\max_{A\in \alph}\|A\|_\infty$, and let $\pi$ be the uniform distribution on the state set of $\tU_m$: $\pi(h) = \frac1{|\tH_m|}$ for every state $h\in \tH_m$. If $n \ge 1$, the distribution vector of the state reached after $n$ random steps in $\tU_m$, starting from $\Id$, namely $\tP_m^n(\Id,\cdot)$, satisfies
\begin{equation}\label{eq:total variation}
\left\| \tP_m^n(\Id,\cdot) - \pi \right\|_{\text{Var}}
\leq \frac12\ \sqrt{|\tH_m|}\ \left(1-\frac1{4k^2\,|\tH_m|^2}\right)^{n}.
\end{equation}
\end{lemma}

\begin{proof}
%
The proof is a combination of results in~\cite{diaconis1991geometric}, on the maximal and minimal eigenvalues of a Markov chain, provided the chain in question is primitive and reversible. The chain $\tU_m$ satisfies these hypotheses, its stationary distribution is the uniform distribution, and its eigenvalues $\beta_i$ ($i \in [0, |\tH_m|-1]$) are real, say $1=\beta_0>\beta_1\geq \beta_2\geq\ldots > \beta_{|\tH_m|-1} > -1$, see \cref{thm: basics on Markov}.

Let $\beta_* = \max\{\beta_1,|\beta_{|\tH_m|-1}|\}$. 
Equation~(1.9) of~\cite[Prop. 3]{diaconis1991geometric} states that
%
$$\left\| \tP_m^n(\Id,\cdot) - \pi \right\|_{\text{Var}} \enspace\le\enspace \frac12\ \sqrt{\frac{1-\pi(\Id)}{\pi(\Id)}}\ \beta_*^n.$$
%
Since 
$\pi(h) = \frac1{|\tH_m|}$ for every $h\in \tH_m$, we have $\frac{1-\pi(\Id)}{\pi(\Id)} = |\tH_m|-1$. It follows that
%
\[
\left\| \tP_m^n(\Id,\cdot) - \pi \right\|_{\text{Var}} \enspace\le\enspace \frac12\ \sqrt{|\tH_m|}\ \beta_*^n.
\]
%
Thus, we only need to prove that $\beta_*\leq 1-\frac1{4k^2\,|\tH_m|^2}$.

We first establish that $\beta_1\leq 1-\frac1{4k^2\,|H_m|^2}$ using Poincaré's inequality, as presented in and with the notation of \cite[Prop. 1]{diaconis1991geometric}. 

For each edge $e$ of $\tU_m$, from state $x$ to state $y$, we let $Q(e) = \tP_m(x,y)\, \pi(x) = \tP_m(y,x)\, \pi(y)$. Hence for every edge $e$ we have
\[
\frac1{4k^2|\tH_m|}\leq Q(e) = \frac1{|\tH_m|}\tP_m(x,y) \leq \frac1{2k|\tH_m|}
\]


For each ordered pair of distinct states $x,y\in \tH_m$, we fix a path $\gamma_{x,y}$ from $x$ to $y$ such that a given edge appears at most once in the path, of length at most $|\tH_m|$.
Such  a path exists since the chain is irreducible.
We define its path length as $|\gamma_{x,y}|_Q = \sum_{e \in \gamma_{x,y}} Q(e)^{-1}$.
Then  $|\gamma_{x,y}|_Q \le 4k^2|\tH_m|^2$.

Poincaré's inequality \cite[Prop. 1]{diaconis1991geometric}  states that
%
$$\beta_1 \le 1 - \frac1\kappa,\enspace\textrm{where}\enspace\kappa = \max_{e\text{ edge}}\sum_{x,y\textrm{ such that }e\in \gamma_{x,y}}|\gamma_{x,y}|_Q \,\pi(x)\,\pi(y).$$ 
%
For all $x,y$, we have $|\gamma_{x,y}|_Q\, \pi(x)\,\pi(y) \le 4k^2$ since $\pi(x)=\pi(y)= \frac1{|\tH_m|}$. Thus $\kappa \le 4k^2|\tH_m|^2$. It follows that $\beta_1\leq 1-\frac1{4k^2\,|\tH_m|^2}$, as announced.

Now we use~\cite[Prop. 2]{diaconis1991geometric} to prove that $\beta_{|\tH_m|-1}\geq \frac1k-1$. For each state $x$, we let $\sigma_x$ be the trajectory consisting of (a single iteration of) the self-loop at $x$. By \cite[Prop. 2]{diaconis1991geometric}, we have
%
$$\beta_{|\tH_m|-1} \ge -1 + \frac2\iota,\enspace\textrm{where}\enspace\iota = \max_{e\text{ edge}}\sum_{x\textrm{ such that }e\in \sigma_x}|\sigma_x|_Q\, \pi(x).$$ 
%
Since each $\sigma_x$ contains only one edge $e$ which
%is an $\epsilon$-step at $x$ and
satisfies $|\sigma_x|_Q = 2k|\tH_m|$, we have $\iota = 2k$ since $\pi(x) = \frac1{|\tH_m|}$ and hence $\beta_{|H_m|-1} \geq -1+\frac1k$. As a result, $|\beta_{|\tH_m|-1}| < \beta_1$ and therefore $\beta_* = \beta_1 \le 1-\frac1{4k\,|\tH_m|^2}$, thus concluding the proof.
%
\end{proof}



%%%%%%%%%%%%%%%%%%%%%%%
\subsection{Proof of \cref{thm: technical theorem}}\label{sec: proof of technnical theorem}

We use the following linear algebraic results. The first one, \cref{lm: matrices of small order}, is a slight generalization of~\cite[Lemma~11]{ParKurlberg2003} to matrices of dimension greater than $2$. \cref{lemma: prime divisor of q(n)} follows from \cref{lm: matrices of small order} and the prime number theorem.

\begin{lemma}\label{lm: matrices of small order}
Let $A\in\GL_d(\mathbb{Z})$ be a matrix of infinite order.  The number of primes $p$ such that $A_p$ has order at most $\mathfrak{o}$ in $\GL_d(\Z/p\Z)$ is $\O(\mathfrak{o}^2)$. 
\end{lemma}

\begin{proof}
%
An immediate induction establishes that, for any $\mathfrak{n}\geq 1$, we have
$\|A^\mathfrak{n}\|_\infty \leq d^{\mathfrak{n}-1}\|A\|_\infty^\mathfrak{n}$. It follows that 
\[
\|A^\mathfrak{n}-\Id\|_\infty \leq d^{\mathfrak{n}-1}\|A\|_\infty^\mathfrak{n} + 1 \leq \left(d\ \|A\|_\infty\right)^\mathfrak{n}.
\]
Since $A$ has infinite order, $A^\mathfrak{n}\neq \Id$ and $\|A^\mathfrak{n}-\Id\|_\infty\neq 0$ for every $\mathfrak{n}\ge 1$.
Moreover, $A_p$ always has finite order in $\GL_d(\Z/p\Z)$. If $A_p$ is of order $\mathfrak{n}$, then $p$ divides $\|A^\mathfrak{n}-\Id\|_\infty$, since every coefficient of $A^\mathfrak{n}-\Id$ is $0$ modulo $p$.

Observe that if an integer $N$ has $x$ distinct prime factors, then $N \ge 2^x$, and hence $x \le \log N$. Let $N = \prod_{\mathfrak{n}=1}^\mathfrak{o}\|A^\mathfrak{n}-\Id\|_\infty$. If $p$ is a prime number such that $A_p$ has order at most $\mathfrak{o}$ in $\GL_d(\Z/p\Z)$, then $p$ must divide $N$. Thus the number of primes $p$ such that $A_p$ has order at most $\mathfrak{o}$ is at most
\[
\log N\leq \log \prod_{\mathfrak{n}=1}^\mathfrak{o} \|A^\mathfrak{n}-\Id\|_\infty
\leq \sum_{\mathfrak{n}=1}^\mathfrak{o} \mathfrak{n}\log \left(d\,\|A\|_\infty\right),
\]
which is $\O(\mathfrak{o}^2)$ when $d$ and $A$ are fixed.
%
\end{proof} 

\begin{corollary}\label{lemma: prime divisor of q(n)}
%
Let $A \in \GL_d(\Z)$ be a matrix of infinite order, and let $q(n)$ be the function defined in \cref{def: def of q}. For each $n$ large enough, $q(n)$ admits a prime factor $p_n$ such that $A_{p_n}$ has  order at least $2\,\log^2 n$ in $\GL_d(\Z/p_n\Z)$.
%
\end{corollary}

\begin{proof}
    %
According to the prime number theorem~\cite{hadamard1896distribution,de1896recherches}, which states that the number of prime numbers less than or equal to $N$ is asymptotically equal to $\frac{N}{\ln N}$, where $\ln$ denote the Napierian logarithm, there exists a positive constant $C$ such that, for $N$ large enough, this number is at least $C\, \frac{N}{\log N}$. Therefore $q(n)$ is a product of at least $C\,\frac {\log^5 n}{5\log \log n}$ different prime numbers.

Since this quantity is asymptotically greater than $4\,\log^4n$, \cref{lm: matrices of small order} establishes that $q(n)$ has a prime factor $p_n$ such that $|\langle A\rangle_{p_n}| > 2\,\log^2 n$.
    %
\end{proof}

Connecting \cref{lemma: prime divisor of q(n)} with \cref{lm:convergence rate}, we get the following result.

\begin{corollary}\label{cor:convergence uniform lazy}
%
Let $q(n)$ be the map defined in \cref{def: def of q}, and let $(h_n)_n$ be a sequence such that $h_n \in \tH_{q(n)}$ for each $n$. 
%
If $H$ is infinite, then the probability that a length $n$ trajectory in $\tU_{q(n)}$ ends in $h_n$ is $\O(\log^{-2}n)$.
%
\end{corollary}


\begin{proof}
%
By a theorem due to Schur \cite{1911:Schur}, a finitely generated subgroup of $\GL_d(\C)$ where each element has finite order must be finite. As a result, since $H$ is infinite, it contains an element $A$ with infinite order. Then by \cref{lemma: prime divisor of q(n)}, for each $n$ large enough, $q(n)$ admits a prime factor $p_n$ such that $A_{p_n}$ is of order at least $2\,\log^2 n$ in $\GL_d(\Z/p_n\Z)$. 

It follows that $|\tH_{p_n}| \ge \frac12\,|H_{p_n}| \ge \frac12\,|\langle A_{p_n}\rangle| \ge \log^2 n$. Moreover, one has $|\tH_{p_n}| \le p_n^{d^2} \le (\log n)^{5d^2}$, since each prime divisor of $q(n)$ is at most equal to $\log^5n$.

It follows from these two inequalities that $\log^2 n \le |\tH_{p_n}| \le p_n^{d^2}$, and hence $p_n\geq (\log n)^{2/d^2}$. Thus, for $n$ sufficiently large, $p_n$ is greater than $2\,\max_{B\in \alph}\|B\|_\infty$, as required to apply \cref{lm:convergence rate}.


For each even integer $n=2\nu$, let $\hat h_n$ be the projection of $h_n$ modulo $p_n$, which is well defined since $p_n$ divides $q(n)$.
If $\hat h_n \notin \tH_{p_n}$, then $\tP_{p_n}^\nu(\Id,\hat h_n) = 0$. If $\hat h_n\in \tH_{p_n}$, by \cref{lm:convergence rate}, and with the notation of that statement, we have
%
\[
\left|\tP_{p_n}^\nu(\Id,\hat h_n) -  \frac1{|\tH_{p_n}|}\right| \le \left\| \tP_{p_n}^\nu(\Id,\cdot) - \pi \right\|_{\text{Var}}
\leq \frac12\ \sqrt{|\tH_{p_n}|}\ \left(1-\frac1{4k^2\,|\tH_{p_n}|^2}\right)^{\nu}.
\]

Since $\log^2n\leq |\tH_{p_n}|\leq (\log n)^{5d_2}$, we have
%
\begin{align*}
%
\tP_{p_n}^\nu(\Id,\hat h_n) &\le \frac1{|\tH_{p_n}|} + \frac12\ \sqrt{|\tH_{p_n}|}\  \left(1-\frac1{4k^2\,|\tH_{p_n}|^2}\right)^{\nu} \\
%
&\le \frac{1}{\log^{2}n} + \frac12\ \sqrt{(\log n)^{5d^2}}\ \exp\left(-\frac{n}{8k^2(\log n)^{10d^2}}\right).
%
\end{align*}
Therefore, $P_{p_n}^n(\Id,\hat h_n) = \tP_{p_n}^\nu(\Id,\hat h_n)$ is $\O(\log^{-2}n)$ if $n$ is even and $\hat h_n\in\tH_{p_n}$. This also holds if
$\hat h_n\notin \tH_{p_n}$, as the corresponding probability is equal to zero.

Now suppose that $n$ is odd, $n=2\nu+1$.  In that case, $P_{p_n}^n(\Id,\hat h_n) = \left(\tP_{p_n}^\nu\times P_{p_n}\right)(\Id,\hat h_n)$. This is equal to $\sum_{k_n \in \hat H_{p_n}}\tP_{p_n}^\nu(\Id, k_n)\, P_{p_n}(k_n,\hat h_n)$. Observe that $P_{p_n}(k_n,\hat h_n) = 0$ unless $k_n = \hat h_n\, A$ for some $A\in \alph$ (which holds for exactly $2k$ values of $k_n$), in which case $P_{p_n}(k_n,\hat h_n) \le \frac 1{2k}$. It follows that, in this case as well, $P_{p_n}^n(\Id,\hat h_n)$ is $\O(\log^{-2}n)$.

Finally, we note that, if $A\in\GL_d(\Z)$ is such that $A_{q(n)}=\Id$, then $A_{p_n}=\Id$ since $p_n$ is a divisor of $q(n)$.
Thus, the probability that a trajectory in $\tU_{q(n)}$ ends in $h_n$ is bounded above by the probability that a trajectory following the same steps in $\tU_{p_n}$ ends in $\hat h_n$, thus concluding the proof.
%
\end{proof}

\cref{cor:convergence uniform lazy} directly implies the proof of 
\cref{thm: technical theorem} by taking each $h_n$ to be $\Id$.
 
\begin{remark}
    It is interesting to note that, in the proof of \cref{cor:convergence uniform lazy}, we are not concerned with the value of the matrix $A$ or the prime $p_n$, nor with how hard it would be to compute them. It is enough, for our purpose, to know that they exist.
\end{remark}

%%%%%%%%%%%%%%%%%%%%%%%
\section{Concluding remarks}\label{sec: concluding}

%%%%%%%%%%%%%%%%%%%%%%%
\subsection{A methodological comment}

In Algorithm $\S$, we compute $\mat(w)$ modulo $q(n)$, the product of all the prime numbers less than or equal to $\log^5n$. The linear average-case complexity is established by the existence of a prime divisor $p_n$ of $q(n)$ such that $H_{p_n} = \langle \Sigma\rangle_{p_n}$ has cardinality at least $\O(\log^2n)$.

Another approach is the following probabilistic algorithm: draw uniformly at random a prime number $p \le \log^7n$ and then modify Algorithm $\S$ to compute $\mat(w)_p$ (instead of $\mat(w)_{q(n)})$.

As in Section~\ref{sec: proof of technnical theorem}, Lemma~\ref{lm: matrices of small order} shows that $\O(\log^4n)$ prime numbers $p'$ are such that $|H_{p'}| \le \log^2n$, so $|H_{p}| \le \log^2n$ with probability $\O\left(\frac{\log^4n}{\frac{\log^7n}{\log\log n}}\right) = \O(\log^{-2}n)$. In particular, this modification of Algorithm $\S$ also works, but a full analysis of its complexity requires evaluating the time needed to draw the prime number $p \le \log^7n$.

One solution is to use Eratosthenes's sieve, as in the proof of Proposition~\ref{prop: t polylogarithmic}, to compute all the primes at most $\log^7n$, and then draw one of those primes uniformly at random. The listing can be done in time $\O(\log n\,\log\log n\, \log\log\log n)$. Another is to use a reject algorithm to draw a number at most equal to $\log^7n$ and test its primality (in polylog time). The expected number of iterations (rejects) is $\O(\log\log n)$.

With either strategy, we get a probabilistic algorithm with the same average-case complexity as our deterministic algorithm, namely $\O(n)$. Therefore, there is no reason to prefer such an algorithm. 


%%%%%%%%%%%%%%%%%%%%%%%
\subsection{Generalization: from $\Z$ to certain subrings of $\C$}\label{sec: generalization}

Let $R$ be a subring of the complex field $\C$, such that $(R,+)$ is a finitely generated group (and hence a free abelian group). In order to compute in $R$, we assume that we are given a basis $(\alpha_1,\dots,\alpha_\delta)$ of $(R,+)$, as well as the expression of the products $\alpha_i\alpha_j$ in that basis. The group $(R,+)$ is isomorphic to $\Z\alpha_1 \oplus \dots \oplus \Z\alpha_\delta$, and its elements are identified with elements of $\Z^\delta$.

The Word Problem is formulated in $\GL_d(R)$ as in $\GL_d(\Z)$.
A fixed finite subset $\Sigma$ of $\GL_d(R)$ is given, along with the inverses $A\inv$ of the matrices $A\in \Sigma$, and $\alph$ denotes the set $\Sigma \cup \Sigma\inv$. Given a word $w\in \alph^*$, one must decide whether the evaluation $\mat(w)$ of $w$ in $\GL_d(R)$ is the identity matrix.

If $m\in \Z$, $m \ge 2$, the set $mR$ is an ideal of $R$. If $A\in R$ or $A \in \GL_d(R)$, we denote by $A_m$ the image of $A$ in $R_m = R/mR$, which is isomorphic to $\bigoplus_i (\Z/m\Z)\, \alpha_i$.

One can then consider the same algorithms $\DC_\Sigma$ and $\S$ (for the same function $q(n)$ defined in \cref{sec: better algorithm}) as for matrices over $\Z$, computing in $R_{q(n)}$ instead of $\Z/q(n)\Z$. The following statement, analogous to \cref{{prop: complexity DC},thm: main theorem}, holds.

\begin{theorem}\label{thm: main for integer rings}
    Let $R$ be a subring of $\C$ such that the additive group $(R,+)$ is finitely generated. The Word Problem $\WP_\Sigma$ in $\GL_d(R)$ has worst-case complexity $\O(n\log^2n)$ if $\langle\Sigma\rangle$ is infinite, and $\O(n)$ if $\langle\Sigma\rangle$ is finite, or if $\Sigma$ consists of upper triangular matrices.
    
    Moreover Algorithm $\S$ solves $\WP_\Sigma$ in $\GL_d(R)$ with linear average-case complexity.
\end{theorem}

\begin{proof}
%
The proof is essentially the same as for matrices over $\Z$. The following is a discussion of the necessary adjustments of the reasoning.

If $x = \sum_{i=1}^\delta x_i\alpha_i$ is an element of $R$ (where each $x_i$ is in $\Z$), we let $\|x\|_\infty = \max_i|x_i|$. If $A\in \GL_d(R)$, we let $\|A\|_\infty = \max\{\|a_{i,j}\|_\infty \mid 1\le i,j\le d\}$. Finally, we let $\beta = \max\{\|\alpha_i\alpha_j\|_\infty \mid 1\le i,j \le \delta\}$. We also let $\ell(x) = \ell(\|x\|_\infty)$ and $\ell(A) = \ell(\|A\|_\infty)$.
%Recall that, for any integer $n$, $\ell(n) = \lceil\log(|n|+1)\rceil+1$ and $|n| \le 2^{\ell(n)}$, see \cref{sec: computing with integers}.
%
The following facts are verified using \cref{prop: basic facts}.
%
\begin{itemize}
    \item If $x_1,\dots, x_k \in R$ with each $\ell(x_i) \le L$, then $\sum_i x_i$ is computed in time $\O(L)$ and $\|\sum_ix_i\|_\infty \le k \max(\|x_i\|_\infty)$.

    \item If $x, y \in R$ and $\ell(x), \ell(y) \le L$, then $xy$ is computed in time $\O(L\log L)$ and $\|xy\|_\infty \le \beta\, \|x\|_\infty \|y\|_\infty$.

    \item If $A, B \in \GL_d(R)$ and $\ell(A), \ell(B) \le L$, then $\|AB\|_\infty \le d\beta\, \|A\|_\infty \|B\|_\infty$, and it is computed in time $\O(L\log L)$.

    \item Suppose now that $y \in R$ is a constant and $x\in R$. Then $\|xy\|_\infty = \O(\|x\|_\infty)$, and $xy$ is computed in time $\O(\ell(x))$. If $B$ is a constant matrix and $A \in \GL_d(R)$, the size of the coefficients of $AB$ is $\O(\|A\|_\infty)$, and $AB$ is computed in time $\O(\ell(A))$.
%
\end{itemize}

The facts listed above, along with \cref{rk: worst case in other rings}, show that the worst case complexity of $\DC_\Sigma$ (on $\GL_d(R)$) is $\O(n\log^2n)$ if $H$ is infinite, and $\O(n)$ if $H$ is finite. The reasoning in \cref{prop: upper triangular} also applies in this case: if $\Sigma$ consists only of upper triangular matrices, then the worst case complexity of $\DC_\Sigma$ is $\O(n)$.

As for matrices over $\Z$, \cref{lemma: ac of DCm} applies and it reduces the proof to proving an analogue of \cref{thm: technical theorem}: verifying that, if $H = \langle\Sigma\rangle$ is infinite, then $\mat(w)_{q(n)} = \Id$ with probability $\O(\log^{-2} n)$ (when $w$ is chosen uniformly at random among length $n$ words).

A minor remark concerns the proof of \cref{lm: matrices of small order} when $\Z$ is replaced by $R$. We note that, if $A$ is a matrix of infinite order with coefficients in $R$, then $\|A^n\|_\infty \le (d\beta)^{n-1}\|A\|_\infty^n$ (instead of $d^{n-1}\|A\|_\infty^n$). This is no obstacle in the rest of the proof of that lemma.

Similarly, the proof of \cref{cor:convergence uniform lazy} carries as in the case of coefficients in $\Z$, since Schur's theorem applies to all finitely generated subgroups of $\GL_d(\C)$, such as $H = \langle\Sigma\rangle$.
%
\end{proof}


%%%%%%%%%%%%%%%%
\subsubsection*{Acknowledgments}
%
The authors gratefully acknowledge conversations with Denis Osin, Aditya Karnataki and Patrick Polo, which helped shape \cref{sec: generalization}.

%%%%%%%%%%%%%%%%%%%
%%%%%%%%%%%%%%%%%%%

\bibliographystyle{alphaurl}
\newcommand{\etalchar}[1]{$^{#1}$}
\begin{thebibliography}{ECH{\etalchar{+}}92}

\bibitem[BBR93]{1993:BabaiBealsRockmore}
Laszlo Babai, Robert Beals, and Daniel Rockmore.
\newblock Deciding finiteness of matrix groups in deterministic polynomial
  time.
\newblock In {\em Proceedings of the 1993 International Symposium on Symbolic
  and Algebraic Computation, ISSAC'93}, pages 117--126. Association for
  Computing Machinery, 1993.

\bibitem[Boo59]{1959:Boone}
William~W. Boone.
\newblock The word problem.
\newblock {\em Ann. Math.}, 70:207--265, 1959.

\bibitem[CLRS22]{cormen2022introduction}
Thomas~H. Cormen, Charles~E. Leiserson, Ronald~L. Rivest, and Clifford Stein.
\newblock {\em Introduction to algorithms}.
\newblock MIT press, 2022.

\bibitem[Deh11]{1911:Dehn}
Max Dehn.
\newblock Über unendliche diskontinuierliche {G}ruppen.
\newblock {\em Math. Annal.}, 71:116--144, 1911.

\bibitem[DF09]{2009:DetinkoFlannery}
Alla Detinko and David Flannery.
\newblock On deciding finiteness of matrix groups.
\newblock {\em J. Symb. Comput.}, 44(8):1037--1043, 2009.

\bibitem[DLVP96]{de1896recherches}
Charles De~La Vallee-Poussin.
\newblock Recherches analytiques sur la th{\'e}orie des nombres premiers.
\newblock {\em Ann. Soc. Sc. Bruxelles}, 1896.

\bibitem[DS91]{diaconis1991geometric}
Persi Diaconis and Daniel Stroock.
\newblock Geometric bounds for eigenvalues of {M}arkov chains.
\newblock {\em Ann. Appl. Probab.}, 1(1):36--61, 1991.

\bibitem[ECH{\etalchar{+}}92]{1992:Epsteinetal}
David B.~A. Epstein, James~W. Cannon, Derek~F. Holt, Silvio V.~F. Levy,
  Michael~S. Paterson, and William~P. Thurston.
\newblock {\em Word processing in groups}.
\newblock Jones and Bartlett Publishers, Boston, MA, 1992.

\bibitem[Had96]{hadamard1896distribution}
Jacques Hadamard.
\newblock Sur la distribution des z{\'e}ros de la fonction $\zeta (s) $ et ses
  cons{\'e}quences arithm{\'e}tiques.
\newblock {\em Bull. Soc. Math. Fr.}, 24:199--220, 1896.

\bibitem[Hir38]{1938:Hirsch}
Kurt~A. Hirsch.
\newblock On infinite soluble groups. {II}.
\newblock {\em Proc. Lond. Math. Soc.}, 44(3):336--344, 1938.

\bibitem[HvdH21]{harvey2021integer}
David Harvey and Joris van~der Hoeven.
\newblock Integer multiplication in time $\mathcal{O}(n\, \log n)$.
\newblock {\em Ann. Math.}, 193(2):563--617, 2021.

\bibitem[Jac78]{1978:Jacob}
Gérard Jacob.
\newblock La finitude des représentations linéaires des semi-groupes est
  décidable.
\newblock {\em J. Algebra}, 52(2):437--459, 1978.

\bibitem[Kur03]{ParKurlberg2003}
Pär Kurlberg.
\newblock On the order of unimodular matrices modulo integers.
\newblock {\em Acta Arith.}, 110(2):141--151, 2003.

\bibitem[LP17]{LevinPeres}
David Levin and Yuval Peres.
\newblock {\em Markov chains and mixing times, Second edition}.
\newblock American Mathematical Society Press, 2017.

\bibitem[LS01]{1977:LyndonSchupp}
Roger~C. Lyndon and Paul~E. Schupp.
\newblock {\em Combinatorial group theory}.
\newblock Classics in Mathematics. Springer-Verlag, Berlin, 2001.
\newblock Reprint of the 1977 edition.

\bibitem[Mai77]{mairson1977some}
Harry~G Mairson.
\newblock Some new upper bounds on the generation of prime numbers.
\newblock {\em Commun. ACM}, 20(9):664--669, 1977.

\bibitem[MKS66]{1966:MagnusKS}
Wilhelm Magnus, Abraham Karrass, and Donald Solitar.
\newblock {\em Combinatorial group theory: {P}resentations of groups in terms
  of generators and relations}.
\newblock Interscience Publishers [John Wiley \& Sons], New York-London-Sydney,
  1966.

\bibitem[Nov55]{1955:Novikov}
Petr~S. Novikov.
\newblock {\em Ob algoritmi\v cesko\u\i\ nerazre\v simosti problemy to\v
  zdestva slov v teorii grupp (On the algorithmic unsolvability of the word
  problem in group theory)}.
\newblock Izdat. Akad. Nauk SSSR, Moscow, 1955.
\newblock Trudy Mat. Inst. Steklov. no. 44.

\bibitem[OS25]{2025:OlshShp}
Alexander Olshanskii and Vladimir Shpilrain.
\newblock Linear average-case complexity of algorithmic problems in groups.
\newblock {\em J. Algebra}, 668:390--419, 2025.

\bibitem[Pri87]{1987:Pritchard}
Paul Pritchard.
\newblock Linear prime-number sieves: a family tree.
\newblock {\em Sci. Comput. Programming}, 9(1):17--35, 1987.

\bibitem[Sch11]{1911:Schur}
Issai Schur.
\newblock Über {G}ruppen periodischer {S}ubstitutionen.
\newblock {\em Sitzungsber. Preuss. Akad. Wiss.}, pages 619--627, 1911.

\bibitem[Sim73]{1973:Simmons}
Harold Simmons.
\newblock The word problem for absolute presentations.
\newblock {\em J. London Math. Soc.}, 6:275--280, 1973.

\bibitem[vzGG13]{vzGathenGerhard}
Joachim von~zur Gathen and Jürgen Gerhard.
\newblock {\em Modern Computer Algebra, Third edition}.
\newblock Cambridge University Press, 2013.

\bibitem[Weh80]{1980:Wehrfritz}
Bertram~A.F. Wehrfritz.
\newblock On finitely generated soluble linear groups.
\newblock {\em Math. Z.}, 170:155--167, 1980.

\end{thebibliography}



\end{document}
